\begin{table*}[!t]
\centering
\scriptsize
\caption{RQ4: Execution errors in CUDA-to-OpenMP kernel translations of \toolname-generated variants by FT-CodeLlama-13B, FT-StarCoderBase-7B, and FT-HPC-Coder-6.7B, ordered by frequency. An execution error is a translation that compiles but does not reproduce the CUDA variant's verdict (PASS or FAIL) from the application's built-in verification: it prints the other verdict or terminates with an error message. \emph{Occurrences} counts the distinct translation mistakes in each category, identified by manually analyzing the failing OpenMP translations from the five mutation runs and three translators; the same mistake made in several translations is counted once. Parentheses give the share of all distinct mistakes. The examples illustrate each category and are not exhaustive. \emph{Apps}: number of applications named in the examples, a lower bound on the applications in which the category occurs. $^\dagger$Categories that we group as mistranslations of CUDA's thread hierarchy or shared memory.}
\label{tab:error_taxonomy_kernel_translation_runtime}
\setlength{\tabcolsep}{3pt}
\begin{tabular}{@{}>{\raggedright\arraybackslash}p{2.1cm}>{\raggedright\arraybackslash}p{4.4cm}>{\raggedright\arraybackslash}p{8.0cm}rr@{}}
\toprule
\textbf{Category} & \textbf{Description} & \textbf{Examples} & \textbf{Occurrences} & \textbf{Apps} \\
\midrule
Floating-Point Precision Loss & Double-precision constants are rewritten as single-precision literals; the rounding difference changes integer results and grows across iterations. & (background-subtract): \texttt{0.92} and \texttt{0.24} become \texttt{0.92f} and \texttt{0.24f}; the \texttt{unsigned char} results differ by~1 and feed the next frame's threshold, so the difference grows to~20, and the application's verification requires an exact match. & 324 (22\%) & 1 \\
\midrule
Multi-dimensional Index Flattening$^\dagger$ & A multi-dimensional thread index space is collapsed into one linear index applied to every dimension, so independent coordinates become identical and threads write to wrong or duplicate locations. & (winograd): \texttt{tile\_i} and \texttt{tile\_j} are both computed as \texttt{omp\_get\_team\_num() * omp\_get\_num\_threads() + omp\_get\_thread\_num()}, collapsing the 2D CUDA grid onto one index; only diagonal tiles (\texttt{tile\_i == tile\_j}) are computed and all others are left unwritten for input (32, 8, 8). & 283 (19\%) & 1 \\
\midrule
Execution Model Assumption Transfer$^\dagger$ & An implicit CUDA execution guarantee (an exact number of threads per block, one block per batch element, or the grid hierarchy) is carried into OpenMP without the guards or restructuring that make it hold. & (aidw): \texttt{tid} is computed as \texttt{team * BLOCK\_SIZE + lid}, but the OpenMP runtime may run fewer threads per team than \texttt{BLOCK\_SIZE}, so output indices are skipped and shared tiles are partly uninitialized. (dense-embedding): under \texttt{target teams distribute} with an inner \texttt{parallel for}, \texttt{batch\_idx} is not isolated per team, so inner threads look up the wrong \texttt{dense\_elem}. (glu): the CUDA thread-index formula is kept but no parallel construct is emitted, so the body executes once and only element~0 is computed. (nlll): each team writes \texttt{*(output + i)} as if every CUDA block had its own output, writing past the single-element output. & 204 (14\%) & 4 \\
\midrule
Intermediate Variable Elimination & A computed intermediate value is dropped and a raw input is used in its place, changing the formula. & (aidw): the normalized ratio \texttt{R\_S0} is dropped and the raw \texttt{r\_obs} is used, so all threshold comparisons and the cosine argument operate at the wrong scale. (glu): the input offset \texttt{xOffset = 2 * split\_dim\_size * N} is dropped and the output stride is used to index the input, so \texttt{x1} and \texttt{x2} read the wrong halves of the input. (winograd): the tile offset is added without the factor~2, or only on the input side, or \texttt{tile\_j\_size} replaces \texttt{MAP\_SIZE} as the row stride, so tiles are read and written at shifted positions. & 198 (13\%) & 3 \\
\midrule
Incorrect API/Library Call & A library math function is replaced by a related but semantically different one, changing the formula. & (lif): \texttt{expm1f} in the voltage update is replaced by \texttt{expf}, dropping the $-1$ term, so every voltage update is off by the full difference between input current and voltage. & 126 (8\%) & 1 \\
\midrule
Shared Memory Scope Mistranslation$^\dagger$ & CUDA \texttt{\_\_shared\_\_} arrays or the synchronization around them are mistranslated: the arrays are privatized per thread, or threads update and read them without a barrier, \texttt{atomic}, or \texttt{reduction}, so the reduction reads incomplete partial results. & (nlll): \texttt{sm\_inputs} and \texttt{acc\_weight} are declared in the \texttt{target teams} region outside \texttt{parallel}, and thread~0's reduction reads incomplete per-thread sums, corrupting \texttt{output} and \texttt{total\_weight} for input (8192, 1000, 10). Other variants declare the arrays inside \texttt{parallel} (private per thread), let every thread zero the whole array without a barrier, or update partial sums in a parallel loop without \texttt{atomic} or \texttt{reduction}. & 111 (7\%) & 1 \\
\midrule
Loop Bound Error$^\dagger$ & A loop that replaces CUDA's thread indexing has the wrong bound, start, or stride, so output elements are skipped or written out of bounds. & (dense-embedding): \texttt{threadIdx.x} is serialized as \texttt{for (tid = 0; tid < numThreads; tid++)}; with \texttt{numThreads = 128} and \texttt{embedding\_dim = 768}, 640 of the 768 columns are never written. (nlll): the partial-sum slot is indexed modulo the block size \texttt{numThreads} instead of \texttt{NLL\_LOSS\_THREADS}, or the CUDA stride is kept with a fixed start of~0, so most frames are skipped. (background-subtract): the loop becomes \texttt{for (i = 1; i <= imgSize; ++i)}, skipping pixel~0 and writing one past the end. & 101 (7\%) & 3 \\
\midrule
Missing Statement Fault & A required statement, typically an in-place state update, is dropped, so stale values corrupt the current output and every later step that reads them. & (lif): \texttt{ref\_time -= dt} is omitted, so the refractory multiplier and the spike-triggered reset use the undecremented value and the spike output is wrong from the first time step for input (1000, 32, 1000); other variants drop the spike reset \texttt{voltage = 0.f}. (flip): \texttt{cur\_indices = rem;} is omitted, so each dimension keeps the quotient instead of the remainder and every source offset is wrong. (ising): the \texttt{is\_black} branch that selects \texttt{joff} and the periodic wrap are dropped, so both sublattices use one neighbor rule and boundary cells read index~$-1$. & 89 (6\%) & 3 \\
\midrule
Incorrect Reduction & A final-result assignment is translated as an accumulation into device memory that is never reset, so the result grows with every kernel launch. & (nlll): \texttt{output} and \texttt{total\_weight} are updated with \texttt{+=} and \texttt{-=} instead of assigned; across the benchmark's 100 launches, \texttt{total\_weight} reaches 100 times the reference value. & 65 (4\%) & 1 \\
\midrule
\textbf{Total} & & & 1{,}501 & 9 \\
\bottomrule
\end{tabular}
\end{table*}
